\documentclass[letter,11pt]{article}
\usepackage[breaklinks]{hyperref}
\hypersetup{
    bookmarks=true,         % show bookmarks bar?
    unicode=false,          % non-Latin characters in Acrobat’s bookmarks
    pdftoolbar=true,        % show Acrobat’s toolbar?
    pdfmenubar=true,        % show Acrobat’s menu?
    pdffitwindow=false,     % window fit to page when opened
    pdfstartview={XYZ null null 1.00},    % disable zoom
    pdftitle={Classwork 8},    % title
    pdfauthor={Richard Zak},     % author
    pdfsubject={UMBC CMSC104 Problem Solving and Computer Programming},   % subject of the document
    pdfkeywords={Computer Science, Programming, Problem Solving, CSEE}, % list of keywords
    pdfnewwindow=true,      % links in new PDF window
    colorlinks=false,       % false: boxed links; true: colored links
    linkcolor=red,          % color of internal links (change box color with linkbordercolor)
    citecolor=green,        % color of links to bibliography
    filecolor=magenta,      % color of file links
    urlcolor=cyan           % color of external links
}
\usepackage{graphicx}
\usepackage{fancyhdr}
\usepackage{multicol}
\pagestyle{fancy}
\usepackage[letterpaper, margin=1in]{geometry}
\geometry{letterpaper}
\usepackage{parskip} % Disable initial indent
\usepackage{color,soul} % Highligher
\usepackage[normalem]{ulem} % Strikethrough with \sout{}
\usepackage{attachfile} % attach files

\usepackage[utf8]{inputenc}
\fancyhf{}
\renewcommand{\headrulewidth}{0pt} % Remove default underline from header package
\rhead{CMSC 104 Section 01: Classwork 8}
%\rhead{}
\lhead{\begin{picture}(0,0) \put(0,-10){\includegraphics[width=1.1cm]{Images/UMBC-vertical}} \end{picture}}
\cfoot{\thepage}
\rfoot{Spring 2024}
\lfoot{CMSC 104 Section 01}
\AtEndDocument{\vfill \footnotesize{Last modified: 08 February 2023}}
\AtEndDocument{\rfoot{Spring 2024}}
\renewcommand\thesubsection{\arabic{subsection}} % Show only subsection numbers, not section.subsection

\begin{document}

\huge
\textbf{Classwork 8: Class Grades Simulator}.
\normalsize
\\ ~~ \\
\textbf{In-class Date: Thursday 20 April} \\
\textbf{Due Date: Wednesday 26 April}

\section*{Objectives}
\paragraph{}Practice writing a program that uses a \texttt{switch} statement, the \texttt{+=} assignment operator, and a \texttt{for} loop.

\section*{Assignment}
\paragraph{}Write a program to simulate class letter grades (A, B, C, D, or F) and output the count of each letter grade to the screen. To do this, you will utilize a for loop that calculates a random letter grade each time through the loop, then utilizes a switch statement to update the appropriate counter variable using the += assignment operator. Make sure the default case warns the user of invalid letter grade. Outside the loop, you will print the final results.

\subsection*{Example Compilation and Execution}
\begin{verbatim}
[rzak1@linux1 cw8]$ gcc -Wall grades.c
[rzak1@linux1 cw8]$ ./a.out
Out of 36 students, here is the class grade breakdown:
A: 10
B: 8
C: 10
D: 3
F: 5
[rzak1@linux1 cw8]$ 
\end{verbatim}

\subsection*{Starter Code}
\paragraph{}Use this code to help you get started.
\begin{verbatim}
/**************************************
** File:  grades.c
** Author: <myName>
** Date:  <todaysDate>
** Section: CMSC104 Section 01
** E-mail: <myEmailAddress>
**
** This file contains the main program for Classwork 8.
** The program randomly generates letter grades (ABCDF),
** counts how many of each there are,
** and prints the results to the screen.
**************************************/

#include<stdio.h>
#include<stdlib.h>
#include<time.h>

#define TOTAL_GRADES 36

int main() {
   // Letter grade counter variables
   // TODO: Initialize each variable appropriately.
   int numA;
   int numB;
   int numC;
   int numD;
   int numF;

   // Variables used in grades loop
   int gradeCounter;
   int totalGrades = TOTAL_GRADES;

   // Variable used to simulate each letter grade
   char letterGrade;

   // Seed the random number generator
   srandom(time(0));

   // Counter-controlled loop to simulate class letter grades
   for (gradeCounter = 0; gradeCounter < totalGrades; gradeCounter++) {
      // Randomly generate a letter grade from ABCDF
      letterGrade = "ABCDF" [random() % 5];

      // TODO: Create a switch statement on letterGrade
      // to update num variables using += assignment operator.
      // NOTE: Don't forget the default case for invalid letterGrade.
   }

   // Print out the final results
   printf("Out of %d grades, here is the class grade breakdown:\n", totalGrades);
   printf("A: %d\n", numA);
   printf("B: %d\n", numB);
   printf("C: %d\n", numC);
   printf("D: %d\n", numD);
   printf("F: %d\n", numF);

   return 0;
}

\end{verbatim}

\subsection*{Notes}
\paragraph{}Notice how the \texttt{for} loop in the program is provided for you, it terminates after the correct number of grades have been simulated, and for the basic assignment, there is no user input.

\paragraph{}Notice how the starting point file uses \texttt{\#define} to define a constant \texttt{TOTAL\_GRADES} that holds the number of grades to simulate. That way, if you want/need to change the total number, you only have to change it in one place.

\section*{Extra Credit}
\paragraph{}Try to do the following embellishments:

\paragraph{Option 1:} After the printout of results, add an if-else block that tells the user if First-Year Intervention (FYI) notifications need to be sent, and if so, how many. To determine if any need to be sent, count how many Ds and Fs there are.

\paragraph{Option 2:} Ask the user for how many grades to simulate. Warn the user if they enter something less than 0. If the user enters 0, print the following to the screen: “Sorry to see you don’t want to use my simulator” with a newline at the end for proper display on screen.

\paragraph{Option 3:} Change the provided for loop to a while loop.

\section*{Grading Rubric}
\begin{itemize}
    \item Header comment: 2 points
    \item Body comments: 3 points
    \item Compiles: 30 points
    \item Proper logic: 55 points
    \item typescript: 10 points
    \item EC1: +5 points
    \item EC2: +5 points
    \item EC3: +5 points
\end{itemize}

\section*{What to Submit}
\paragraph{}Use the \texttt{script} command to record yourself compiling and running your program 3 times. (Do not record yourself editing your program!) Exit from script. Submit your program and the typescript file.
\begin{verbatim}
[rzak1@linux1 cw8]$ submit cmsc104_rzak1 cw8 grades.c typescript
\end{verbatim}

\subsection*{Verify}
\paragraph{}Make sure you submitted the assignment correctly.
\begin{verbatim}
[rzak1@linux1 cw8]$ submitls cmsc104_rzak1 cw8
\end{verbatim}
\end{document}